\setchapterimage{figure/Chapter4封面图}
\setchapterintro{}
\chapter{常系数线性方程}\label{chap:常系数线性方程}

考虑一般的 $n$ 阶常系数线性微分方程
\[
	\dfrac{\dd^n x}{\dd t^n}
	+a_1\dfrac{\dd^{n-1}x}{\dd t^{n-1}}
	+\cdots
	+a_{n-1}\dfrac{\dd x}{\dd t}
	+a_nx=f(t),
\]
其中 $a_1,a_2,\ldots,a_n$ 都是常数。引入微分算子
\[
	D=\dfrac{\dd}{\dd t},
	\qquad
	D^n=\dfrac{\dd^n}{\dd t^n},
\]
则上述方程可以简写为
\[
	P(D)x=f(t),
\]
其中
\[
	P(D)=D^n+a_1D^{n-1}+\cdots+a_{n-1}D+a_n
\]
是关于 $D$ 的 $n$ 次多项式。

算子 $D$ 可以看作从连续可微函数空间 $C^1(J)$ 到连续函数空间
$C^0(J)$ 的映射，其中 $J$ 是函数的定义域。$C^1(J)$ 是
$C^0(J)$ 的子空间，二者都是线性空间，而 $D$ 是其上的线性算子。

由第三章可知，非齐次线性方程的通解由一个特解与相应齐次方程通解之和构成。
因此，首先研究齐次方程
\[
	\dfrac{\dd^n x}{\dd t^n}
	+a_1\dfrac{\dd^{n-1}x}{\dd t^{n-1}}
	+\cdots
	+a_{n-1}\dfrac{\dd x}{\dd t}
	+a_nx=0.
\]
一旦确定了这个齐次方程的通解，原非齐次方程的特解原则上可由常数变易法得到。

\section{齐次问题}\label{sec:常系数线性方程的齐次问题}

设齐次方程具有形如
\[
	x(t)=\mathrm{e}^{\lambda t}
\]
的解，其中 $\lambda$ 为待定常数。代入齐次方程可得
\[
	P(D)\mathrm{e}^{\lambda t}=P(\lambda)\mathrm{e}^{\lambda t}=0,
\]
从而问题化为多项式方程
\[
	P(\lambda)
	=\lambda^n+a_1\lambda^{n-1}+\cdots+a_{n-1}\lambda+a_n=0
\]
的求根问题。多项式 $P(\lambda)$ 称为\textbf{特征多项式}，
$P(\lambda)=0$ 称为\textbf{特征方程}，它的根称为\textbf{特征根}；
相应的指数函数解称为\textbf{特征解}。

\begin{theorem}\label{thm:互异特征根的基本解组}
	若特征方程有 $n$ 个互异的根
	$\lambda_1,\lambda_2,\ldots,\lambda_n$，则齐次方程有基本解组
	\[
		\mathrm{e}^{\lambda_1t},\mathrm{e}^{\lambda_2t},\ldots,\mathrm{e}^{\lambda_nt}.
	\]
\end{theorem}

\begin{proof}
	这些函数显然都是齐次方程的解。它们的 Wronski 行列式为
	\[
		\begin{aligned}
		W(t)
		&=
		\begin{vmatrix}
			1&1&\cdots&1\\
			\lambda_1&\lambda_2&\cdots&\lambda_n\\
			\lambda_1^2&\lambda_2^2&\cdots&\lambda_n^2\\
			\vdots&\vdots&&\vdots\\
			\lambda_1^{n-1}&\lambda_2^{n-1}&\cdots&\lambda_n^{n-1}
		\end{vmatrix}
		\mathrm{e}^{(\lambda_1+\lambda_2+\cdots+\lambda_n)t}\\
		&=\prod_{1\leq j<k\leq n}(\lambda_k-\lambda_j)\mathrm{e}^{-a_1t}\neq0.
		\end{aligned}
	\]
	这里使用了 Vandermonde 行列式以及根与系数的关系
	$\lambda_1+\cdots+\lambda_n=-a_1$。故这些解线性无关，因而构成基本解组。
\end{proof}

\begin{theorem}\label{thm:重特征根的基本解组}
	设特征方程只有 $r$ 个互异的根
	$\lambda_1,\lambda_2,\ldots,\lambda_r$，其重数分别为
	$n_1,n_2,\ldots,n_r$，其中
	\[
		n_1+n_2+\cdots+n_r=n.
	\]
	则对每个 $i=1,2,\ldots,r$，函数
	\[
		\mathrm{e}^{\lambda_it},
		t\mathrm{e}^{\lambda_it},
		\ldots,
		t^{n_i-1}\mathrm{e}^{\lambda_it}
	\]
	都是齐次方程的解；将各组函数合在一起，便得到该方程的一个基本解组。
\end{theorem}

\begin{proof}
	证明分为两步。

	先固定 $i$。对任意非负整数 $l$，由 Leibniz 公式或算子恒等式
	$D^m(\mathrm{e}^{\lambda t}u)=\mathrm{e}^{\lambda t}(D+\lambda)^mu$，有
	\[
		P(D)\bigl(t^l\mathrm{e}^{\lambda t}\bigr)
		=\mathrm{e}^{\lambda t}P(D+\lambda)t^l
		=\mathrm{e}^{\lambda t}
		\sum_{k=0}^{l}\binom{l}{k}P^{(k)}(\lambda)t^{l-k}.
	\]
	若 $\lambda_i$ 是 $n_i$ 重根，则
	\[
		P(\lambda_i)=P'(\lambda_i)=\cdots
		=P^{(n_i-1)}(\lambda_i)=0.
	\]
	因此当 $l=0,1,\ldots,n_i-1$ 时，
	$t^l\mathrm{e}^{\lambda_it}$ 都是齐次方程的解。若
	\[
		\sum_{k=0}^{n_i-1}C_kt^k\mathrm{e}^{\lambda_it}\equiv0,
	\]
	由于 $\mathrm{e}^{\lambda_it}\neq0$，便有
	$\sum\limits_{k=0}^{n_i-1}C_kt^k\equiv0$，从而所有 $C_k$ 均为零。
	故同一特征根对应的这一组解线性无关。

	再证明不同特征根所对应的各组解合在一起仍然线性无关。
	若存在恒等式
	\[
		\mathrm{e}^{\lambda_1t}\phi_1(t)
		+\mathrm{e}^{\lambda_2t}\phi_2(t)
		+\cdots
		+\mathrm{e}^{\lambda_rt}\phi_r(t)\equiv0,
	\]
	其中 $\deg\phi_i<n_i$，则除以 $\mathrm{e}^{\lambda_1t}$ 
	\[
	\phi_1(t)
	+\mathrm{e}^{\lambda_2t-\lambda_1t}\phi_2(t)
	+\cdots
	+\mathrm{e}^{\lambda_rt-\lambda_1t}\phi_r(t)\equiv0,\]
	求一次导数
	\[	\phi_1'(t)
	+\mathrm{e}^{\lambda_2t-\lambda_1t}[\phi_2'(t)+(\lambda_2-\lambda_1)\phi_2(t)]
	+\cdots
	+\mathrm{e}^{\lambda_rt-\lambda_1t}[\phi_r'(t)+(\lambda_r-\lambda_1)\phi_r(t)]\equiv0,\]
	后面的项次数不变，连续作用
	$n_1$ 次微分算子。第一项被消去，其余各项仍是非零指数函数与同次数多项式的乘积。
	再依次消去第二项、第三项，最终得到
	\[
		\mathrm{e}^{\mu t}\pi(t)\equiv0,
	\]
	其中 $\pi(t)$ 是不恒为零的多项式。这与
	$\mathrm{e}^{\mu t}\neq0$ 矛盾。因此所有 $\phi_i$ 都只能恒为零，各组解合在一起线性无关。
	这些解的总数为 $n_1+\cdots+n_r=n$，故它们构成齐次方程的基本解组。
\end{proof}

\noindent\textbf{补充：利用 Jordan 块构造重根对应的解}\quad
固定特征根 $\lambda_i$，记 $m=n_i$。由于
\[
	P(z)=(z-\lambda_i)^mQ_i(z),
\]
只要构造出 $(D-\lambda_i)^mx=0$ 的解，就得到原齐次方程的解。
令
\[
	y_1=x,\qquad y_2=(D-\lambda_i)x,\qquad\ldots,\qquad
	y_m=(D-\lambda_i)^{m-1}x,
\]
则方程 $(D-\lambda_i)^mx=0$ 等价于一阶方程组
\[
	\vect{y}'=J_i\vect{y},\qquad
	J_i=\lambda_iI_m+N_m,\qquad
	N_m=
	\begin{pmatrix}
		0&1&0&\cdots&0\\
		0&0&1&\ddots&\vdots\\
		\vdots&\ddots&\ddots&\ddots&0\\
		0&\cdots&0&0&1\\
		0&\cdots&\cdots&0&0
	\end{pmatrix}_{m\times m}.
\]
这里 $J_i$ 是一个 $m$ 阶 Jordan 块；当 $m=1$ 时，取 $N_1=(0)$。
作代换 $\vect{y}=\mathrm{e}^{\lambda_it}\vect{u}$，方程组化为
\[
	\vect{u}'=N_m\vect{u},
	\qquad
	u_1'=u_2,\quad\ldots,\quad u_{m-1}'=u_m,\quad u_m'=0.
\]
从最后一个方程开始逐次积分，得到
\[
	u_m=c_m,\qquad u_{m-1}=c_{m-1}+c_mt,\qquad\ldots,\qquad
	u_1=c_1+c_2t+\cdots+c_m\dfrac{t^{m-1}}{(m-1)!}.
\]
因此，取第一分量便得到
\[
	x(t)=\mathrm{e}^{\lambda_it}
	\left(c_1+c_2t+\cdots+c_m\dfrac{t^{m-1}}{(m-1)!}\right).
\]
分别令 $(c_1,\ldots,c_m)$ 为各个标准基向量，就构造出
\[
	\mathrm{e}^{\lambda_it},\quad
	t\mathrm{e}^{\lambda_it},\quad\ldots,\quad
	\dfrac{t^{m-1}}{(m-1)!}\mathrm{e}^{\lambda_it}.
\]
将各函数乘以相应的非零常数 $0!,1!,\ldots,(m-1)!$，即得定理中的
$\mathrm{e}^{\lambda_it},t\mathrm{e}^{\lambda_it},\ldots,
t^{n_i-1}\mathrm{e}^{\lambda_it}$。
上述构造也可写成矩阵形式
\[
	\vect{y}(t)=\mathrm{e}^{\lambda_it}
	\left(I_m+tN_m+\cdots+\dfrac{t^{m-1}}{(m-1)!}N_m^{m-1}\right)\vect{c}.
\]
由于 $N_m^m=0$，括号中的有限和就是 $\mathrm{e}^{tN_m}$，
从而整个矩阵为 $\mathrm{e}^{tJ_i}$；矩阵指数的定义及一般性质将在
第~\ref{sec:常系数线性方程组}~节中讨论。

上述结论先是在复数域中讨论的。对实系数方程，成对出现的共轭复根可以转化为实函数解。

\begin{corollary}\label{cor:实系数常系数方程的基本解组}
	设实系数齐次线性方程具有 $r$ 个互异的实特征根
	$\lambda_1,\lambda_2,\ldots,\lambda_r$，其重数分别为
	$n_1,n_2,\ldots,n_r$；又具有 $l$ 对互异的复特征根
	\[
		\alpha_1\pm\mathrm{i}\beta_1,
		\alpha_2\pm\mathrm{i}\beta_2,
		\ldots,
		\alpha_l\pm\mathrm{i}\beta_l,
	\]
	其重数分别为 $m_1,m_2,\ldots,m_l$，并且
	\[
		\sum_{k=1}^{r}n_k+2\sum_{j=1}^{l}m_j=n.
	\]
	则该方程的下列实函数组成一个基本解组：
	\[
		\begin{gathered}
			t^q\mathrm{e}^{\lambda_kt},
			\qquad q=0,1,\ldots,n_k-1,
			\quad k=1,2,\ldots,r,\\
			t^q\mathrm{e}^{\alpha_jt}\cos\beta_jt,
			\qquad q=0,1,\ldots,m_j-1,
			\quad j=1,2,\ldots,l,\\
			t^q\mathrm{e}^{\alpha_jt}\sin\beta_jt,
			\qquad q=0,1,\ldots,m_j-1,
			\quad j=1,2,\ldots,l.
		\end{gathered}
	\]
\end{corollary}

\begin{proof}
	由命题~\ref{thm:实系数齐次方程组的复值解}，实系数齐次线性方程的复值解
	$x(t)=u(t)+\mathrm{i} v(t)$ 的实部 $u(t)$ 与虚部 $v(t)$ 仍是该方程的解。
	将定理~\ref{thm:重特征根的基本解组} 中共轭复根对应的解取实部和虚部，
	便得到上述实函数。原来的 $n$ 个复值解线性无关，且它们都可由上述
	$n$ 个实函数作常系数线性组合表示，因此这些实函数也线性无关，从而构成基本解组。
\end{proof}

\section{非齐次问题}\label{sec:常系数线性方程的非齐次问题}

对一般问题来说，利用常数变易法求非齐次方程的特解，计算可能相当复杂。
针对几类特殊而常见的函数，可以采用更简便的\textbf{算子解法}。

考虑非齐次线性方程
\[
	P(D)x=f(t).
\]
将 $P(D)$ 的逆算子形式地记为 $\dfrac{1}{P(D)}$，则可将方程的一个特解写成
\[
	x(t)=\dfrac{1}{P(D)}f(t).
\]
先说明逆算子记号的含义。例如，$\dfrac{1}{D}$ 作用于可积函数 $f(t)$ 时，
\[
	\dfrac{1}{D}f(t)=\int f(t)\,\dd t.
\]
由于这里只需求任意一个特解，可以在每次积分时选取一个方便的原函数，
而不计积分常数。以下算子等式都在这个意义下理解。

逆算子 $\dfrac{1}{P(D)}$ 具有下列基本性质。

\noindent\textbf{性质 1}\quad
\[
	\dfrac{1}{D^n}f(t)
	=\underbrace{\int\cdots\int}_{n\text{ 次}}f(t)\,(\dd t)^n,
\]
即作 $n$ 次积分。

\noindent\textbf{性质 2}\quad 逆算子的作用是线性的，即
\[
	\dfrac{1}{P(D)}\bigl(\alpha f_1(t)+\beta f_2(t)\bigr)
	=\alpha\dfrac{1}{P(D)}f_1(t)
	+\beta\dfrac{1}{P(D)}f_2(t).
\]

\noindent\textbf{性质 3}\quad 若 $P(D)=P_1(D)P_2(D)$，则
\[
	\dfrac{1}{P(D)}
	=\dfrac{1}{P_1(D)}\dfrac{1}{P_2(D)}
	=\dfrac{1}{P_2(D)}\dfrac{1}{P_1(D)}.
\]

这些性质均可由微分算子 $P(D)$ 的线性及可交换性直接验证。
下面给出三种重要的计算公式。

\begin{proposition}\label{thm:逆算子的三种计算方法}
	逆算子 $\dfrac{1}{P(D)}$ 有以下计算方法。

	\textup{i）解析展开法（解析相除法）}\quad
	设 $f_k(t)$ 是 $k$ 次多项式。若 $\dfrac{1}{P(x)}$ 在 $x=0$ 处解析，且
	\[
		\dfrac{1}{P(x)}=Q_k(x)+H_k(x),
	\]
	其中 $Q_k(x)$ 是 $k$ 次多项式，$H_k(x)$ 只含 $k+1$ 次及以上的项，则
	\[
		\dfrac{1}{P(D)}f_k(t)=Q_k(D)f_k(t).
	\]

	\textup{ii）代换法}\quad 若 $P(\lambda)\neq0$，则
	\[
		\dfrac{1}{P(D)}\mathrm{e}^{\lambda t}
		=\dfrac{1}{P(\lambda)}\mathrm{e}^{\lambda t}.
	\]

	\textup{iii）二项式法}\quad 对任意适当的函数 $v(t)$，有
	\[
		\dfrac{1}{P(D)}\bigl(\mathrm{e}^{\lambda t}v(t)\bigr)
		=\mathrm{e}^{\lambda t}\dfrac{1}{P(D+\lambda)}v(t).
	\]
\end{proposition}

\begin{proof}
	对解析展开法，由于 $f_k(t)$ 是 $k$ 次多项式，故
	$D^{k+i}f_k(t)=0$（$i\geq1$）。因此 $H_k(D)f_k(t)=0$，从而只留下
	$Q_k(D)f_k(t)$。

	对代换法，由
	\[
		P(D)\mathrm{e}^{\lambda t}=P(\lambda)\mathrm{e}^{\lambda t}
	\]
	即可得到结论。

	对二项式法，由 Leibniz 公式，对任意正整数 $m$ 有
	\[
		\begin{aligned}
		D^m\bigl(\mathrm{e}^{\lambda t}v(t)\bigr)
		&=\sum_{j=0}^{m}\binom{m}{j}
		D^j\mathrm{e}^{\lambda t}\,D^{m-j}v(t)\\
		&=\mathrm{e}^{\lambda t}
		\sum_{j=0}^{m}\binom{m}{j}\lambda^jD^{m-j}v(t)\\
		&=\mathrm{e}^{\lambda t}(D+\lambda)^mv(t).
		\end{aligned}
	\]
	于是
	\[
		P(D)\bigl(\mathrm{e}^{\lambda t}v(t)\bigr)
		=\mathrm{e}^{\lambda t}P(D+\lambda)v(t).
	\]
	将 $v(t)$ 换成 $\dfrac{1}{P(D+\lambda)}v(t)$，便得
	\[
		P(D)\left[
		\mathrm{e}^{\lambda t}\dfrac{1}{P(D+\lambda)}v(t)
		\right]
		=\mathrm{e}^{\lambda t}v(t),
	\]
	即得二项式法公式。
\end{proof}

\noindent\textbf{注 4.1}\quad
解析展开法最直接、最常用的一个情形是
\[
	\dfrac{1}{1-D}f_k(t)
	=\bigl(1+D+D^2+\cdots+D^k\bigr)f_k(t).
\]

\noindent\textbf{注 4.2}\quad
对只含 $D^2$ 的多项式，可由代换法得到
\[
	\dfrac{1}{P(D^2)}
	\begin{Bmatrix}
		\cos at\\
		\sin at
	\end{Bmatrix}
	=\dfrac{1}{P(-a^2)}
	\begin{Bmatrix}
		\cos at\\
		\sin at
	\end{Bmatrix},
	\qquad P(-a^2)\neq0.
\]
事实上，对辅助函数 $\mathrm{e}^{\mathrm{i}at}$ 应用代换法，有
\[
	\dfrac{1}{P(D^2)}\mathrm{e}^{\mathrm{i}at}
	=\dfrac{1}{P\bigl((\mathrm{i}a)^2\bigr)}\mathrm{e}^{\mathrm{i}at}
	=\dfrac{1}{P(-a^2)}\mathrm{e}^{\mathrm{i}at}.
\]
分别取实部和虚部即得上述公式。

\noindent\textbf{注 4.3}\quad
当 $P(\lambda)=0$ 时不能直接使用代换法，但仍可使用二项式法：
\[
	\dfrac{1}{P(D)}\mathrm{e}^{\lambda t}
	=\mathrm{e}^{\lambda t}\dfrac{1}{P(D+\lambda)}\,1.
\]

\section{常系数线性方程组}\label{sec:常系数线性方程组}

考虑常系数齐次线性方程组
\[
	\dfrac{\dd}{\dd t}
	\begin{pmatrix}
		x_1\\
		\vdots\\
		x_n
	\end{pmatrix}
	=
	\begin{pmatrix}
		a_{11}&\cdots&a_{1n}\\
		\vdots&&\vdots\\
		a_{n1}&\cdots&a_{nn}
	\end{pmatrix}
	\begin{pmatrix}
		x_1\\
		\vdots\\
		x_n
	\end{pmatrix}.
\]
由于下面需要使用特征值，先在复数域上讨论。令
\[
	\vect{x}=(x_1,x_2,\ldots,x_n)^{\mathsf T},
	\qquad
	A=(a_{ij})_{n\times n},
\]
则方程组可简写为
\[
	\vect{x}'=A\vect{x}.
\]
设它具有形如
\[
	\vect{x}(t)=\vect{c}\,\mathrm{e}^{\lambda t}
\]
的特解，其中 $\lambda\in\C$，非零向量
$\vect{c}=(c_1,c_2,\ldots,c_n)^{\mathsf T}\in\C^n$。代入方程组得到
\[
	(\lambda I-A)\vect{c}=\vect{0}.
\]
因此，问题化为求矩阵 $A$ 的特征值及相应特征向量。方程
\[
	\det(\lambda I-A)=0
\]
称为矩阵 $A$ 的特征方程。

\begin{proposition}\label{thm:互异特征值常系数方程组的基本解组}
	若矩阵 $A$ 有 $n$ 个互异的特征值
	$\lambda_1,\lambda_2,\ldots,\lambda_n$，相应的特征向量为
	$\vect{c}^{(1)},\vect{c}^{(2)},\ldots,\vect{c}^{(n)}$，则方程组
	$\vect{x}'=A\vect{x}$ 有基本解组
	\[
		\vect{c}^{(1)}\mathrm{e}^{\lambda_1t},
		\vect{c}^{(2)}\mathrm{e}^{\lambda_2t},
		\ldots,
		\vect{c}^{(n)}\mathrm{e}^{\lambda_nt}.
	\]
\end{proposition}

\begin{proof}
	每个函数都是方程组的解。互异特征值所对应的特征向量线性无关，故在
	$t=0$ 时，上述 $n$ 个解向量组成的矩阵可逆。由基本解组的判别条件，
	它们构成方程组的基本解组。
\end{proof}

当 $A$ 不可对角化时，需要借助矩阵指数函数与 Jordan 标准形处理。

\subsection*{一、矩阵指数函数}

记全体 $n$ 阶复方阵组成的线性空间为 $\mathscr{M}(n)$，并定义矩阵范数
\[
	\norm{A}=\sum_{i,j=1}^{n}\abs{a_{ij}},
	\qquad A=(a_{ij})\in\mathscr{M}(n).
\]

\begin{lemma}\label{lem:矩阵指数级数的收敛性}
	空间 $\mathscr{M}(n)$ 关于上述范数是完备的，并且对任意
	$A,B\in\mathscr{M}(n)$，有
	\[
		\norm{AB}\leq\norm{A}\norm{B}.
	\]
	特别地，当 $k\geq1$ 时，$\norm{A^k}\leq\norm{A}^k$。此外，矩阵幂级数
	\[
		I+A+\dfrac{1}{2!}A^2+\cdots+\dfrac{1}{k!}A^k+\cdots
	\]
	绝对收敛。
\end{lemma}

\begin{proof}
	完备性由有限维线性空间的完备性得到。矩阵乘法的分量估计直接给出
	$\norm{AB}\leq\norm{A}\norm{B}$，进而归纳得到
	$\norm{A^k}\leq\norm{A}^k$。令 $M=\norm{A}$，则
	\[
		\sum_{k=0}^{\infty}\norm{\dfrac{A^k}{k!}}
		\leq\sum_{k=0}^{\infty}\dfrac{M^k}{k!}
		=\mathrm{e}^{M}<\infty,
	\]
	故矩阵幂级数绝对收敛。
\end{proof}

\begin{definition}\label{def:矩阵指数函数}
	对任意 $A\in\mathscr{M}(n)$，定义
	\[
		\mathrm{e}^{A}=\sum_{k=0}^{\infty}\dfrac{A^k}{k!},
	\]
	称为矩阵 $A$ 的\textbf{指数函数}。
\end{definition}

由幂级数运算可得矩阵指数函数的基本性质：

\begin{enumerate}
	\item 若 $AB=BA$，则
	\[
		\mathrm{e}^{A+B}
		=\mathrm{e}^{A}\mathrm{e}^{B}
		=\mathrm{e}^{B}\mathrm{e}^{A}.
	\]
	\item $\bigl(\mathrm{e}^{A}\bigr)^{-1}=\mathrm{e}^{-A}$。
	\item 若 $P$ 可逆，则
	\[
		\mathrm{e}^{PAP^{-1}}=P\mathrm{e}^{A}P^{-1}.
	\]
\end{enumerate}

利用矩阵指数函数，可以简洁地表述常系数线性方程组的基本结果。

\begin{theorem}\label{thm:常系数线性方程组的矩阵指数解}
	对齐次方程组 $\vect{x}'=A\vect{x}$，矩阵
	\[
		X(t)=\mathrm{e}^{tA}
	\]
	是其标准解矩阵，即 $X(0)=I$。非齐次方程组
	\[
		\vect{x}'=A\vect{x}+\vect{f}(t)
	\]
	的通解为
	\[
		\vect{x}(t)
		=\mathrm{e}^{tA}\vect{c}
		+\int_{t_0}^{t}\mathrm{e}^{(t-s)A}\vect{f}(s)\,\dd s,
	\]
	其中 $\vect{c}$ 为任意常向量。
\end{theorem}

\begin{proof}
	在任意有限区间上，$\mathrm{e}^{tA}$ 的幂级数及其逐项求导级数一致收敛，故
	\[
		\dfrac{\dd}{\dd t}\mathrm{e}^{tA}
		=A+tA^2+\dfrac{t^2}{2!}A^3+\cdots
		=A\mathrm{e}^{tA}.
	\]
	又 $\mathrm{e}^{0A}=I$，所以 $\mathrm{e}^{tA}$ 是标准解矩阵。
	非齐次方程的通解公式则由常数变易公式直接得到。
\end{proof}

\subsection*{二、标准解矩阵的初等表达}

由定义可知，标准解矩阵 $\mathrm{e}^{tA}$ 是一个无穷和。下面利用 Jordan 标准形，
把它写成初等函数的有限和。

存在可逆矩阵 $P$，使得
\[
	P^{-1}AP
	=J
	=\operatorname{diag}(J_1,J_2,\ldots,J_s),
\]
其中 $J_i$ 是 $n_i$ 阶 Jordan 块：
\[
	J_i=
	\begin{pmatrix}
		\lambda_i&1&0&\cdots&0\\
		0&\lambda_i&1&\ddots&\vdots\\
		\vdots&\ddots&\ddots&\ddots&0\\
		0&\cdots&0&\lambda_i&1\\
		0&\cdots&\cdots&0&\lambda_i
	\end{pmatrix}_{n_i\times n_i},
	\qquad i=1,2,\ldots,s,
\]
且 $n_1+n_2+\cdots+n_s=n$。作变量代换 $\vect{x}=P\vect{y}$，
原齐次方程组化为
\[
	\vect{y}'=J\vect{y}.
\]
由矩阵指数函数的相似不变性，
\[
	\mathrm{e}^{tA}=P\mathrm{e}^{tJ}P^{-1},
	\qquad
	\mathrm{e}^{tA}P=P\mathrm{e}^{tJ}
\]
$P\mathrm{e}^{tJ}$也是原方程组的一个基本解矩阵。因此只需计算各 Jordan 块的指数函数。

写成
\[
	J_i=\lambda_iI+Z_i,
	\qquad
	Z_i=
	\begin{pmatrix}
		0&1&0&\cdots&0\\
		0&0&1&\ddots&\vdots\\
		\vdots&\ddots&\ddots&\ddots&0\\
		0&\cdots&0&0&1\\
		0&\cdots&\cdots&0&0
	\end{pmatrix}_{n_i\times n_i}.
\]
矩阵 $Z_i$ 是幂零矩阵，$Z_i^{n_i}=0$，从而
\[
	\mathrm{e}^{tZ_i}
	=I+tZ_i+\dfrac{t^2}{2!}Z_i^2+\cdots
	+\dfrac{t^{n_i-1}}{(n_i-1)!}Z_i^{n_i-1}.
\]
由于 $\lambda_iI$ 与 $Z_i$ 可交换，
\[
	\mathrm{e}^{tJ_i}
	=\mathrm{e}^{\lambda_it}\mathrm{e}^{tZ_i}
	=\mathrm{e}^{\lambda_it}
	\begin{pmatrix}
		1&t&\dfrac{t^2}{2!}&\cdots&\dfrac{t^{n_i-1}}{(n_i-1)!}\\
		0&1&t&\cdots&\dfrac{t^{n_i-2}}{(n_i-2)!}\\
		\vdots&\ddots&\ddots&\ddots&\vdots\\
		0&\cdots&0&1&t\\
		0&\cdots&\cdots&0&1
	\end{pmatrix}.
\]

将 $P$ 按照 Jordan 块分为列块
\[
	P=(P_1,P_2,\ldots,P_s),
	\qquad
	P_i=(\vect{p}_1,\vect{p}_2,\ldots,\vect{p}_{n_i}),
\]
则
\[
	P\mathrm{e}^{tJ}
	=\bigl(P_1\mathrm{e}^{tJ_1},
	P_2\mathrm{e}^{tJ_2},\ldots,
	P_s\mathrm{e}^{tJ_s}\bigr).
\]
其中
\[
	P_i\mathrm{e}^{tJ_i}
	=\mathrm{e}^{\lambda_it}
	\left(
	\vect{p}_1,
	t\vect{p}_1+\vect{p}_2,
	\ldots,
	\dfrac{t^{n_i-1}}{(n_i-1)!}\vect{p}_1
	+\cdots+t\vect{p}_{n_i-1}+\vect{p}_{n_i}
	\right).
\]
因此，与特征值 $\lambda_i$ 对应的每个解向量都具有形式
\[
	\vect{x}(t)
	=\mathrm{e}^{\lambda_it}
	\left[
		\vect{c}_0+\dfrac{t}{1!}\vect{c}_1+\cdots
		+\dfrac{t^{n_i-1}}{(n_i-1)!}\vect{c}_{n_i-1}
	\right].
\]
这就是基本解组的有限和表达。

\subsection*{三、重特征根情形的结论}

设矩阵 $A$ 有互异的特征值
$\lambda_1,\lambda_2,\ldots,\lambda_s$，其代数重数分别为
$n_1,n_2,\ldots,n_s$，且 $n_1+n_2+\cdots+n_s=n$。

\begin{lemma}\label{lem:广义特征子空间分解}
	令
	\[
		V_i
		=\left\{
		\vect{c}\in\C^n:
		(A-\lambda_iI)^{n_i}\vect{c}=\vect{0}
		\right\}.
	\]
	则 $V_i$ 是 $n_i$ 维的 $A$-不变子空间，并且
	\[
		\C^n=V_1\oplus V_2\oplus\cdots\oplus V_s.
	\]
\end{lemma}
\begin{lemma}[待定系数的确定]\label{lem:向量多项式解的递推条件}
	设 $\lambda_i$ 是 $A$ 的 $n_i$ 重特征值，则向量函数
	\[
		\vect{x}(t)
		=\mathrm{e}^{\lambda_it}
		\left[
			\vect{c}_0+\dfrac{t}{1!}\vect{c}_1+\cdots
			+\dfrac{t^{n_i-1}}{(n_i-1)!}\vect{c}_{n_i-1}
		\right]
	\]
	是齐次方程组$\vect{x}'=A\vect{x}$的非零解，当且仅当系数向量满足
	\[
		\begin{cases}
			(A-\lambda_iI)^{n_i}\vect{c}_0=\vect{0},
			\qquad \vect{c}_0\neq\vect{0},\\
			\vect{c}_1=(A-\lambda_iI)\vect{c}_0,\\
			\vect{c}_2=(A-\lambda_iI)\vect{c}_1,\\
			\hspace{2.5cm}\vdots\\
			\vect{c}_{n_i-1}
			=(A-\lambda_iI)\vect{c}_{n_i-2}.
		\end{cases}
	\]
\end{lemma}

\begin{proof}
	将上述待定式代入 $\vect{x}'=A\vect{x}$，消去公因子
	$\mathrm{e}^{\lambda_it}$ 后比较 $t$ 的同次幂，得到
	\[
		(A-\lambda_iI)\vect{c}_0=\vect{c}_1,
		\quad\ldots,\quad
		(A-\lambda_iI)\vect{c}_{n_i-2}=\vect{c}_{n_i-1},
		\quad
		(A-\lambda_iI)\vect{c}_{n_i-1}=\vect{0}.
	\]
	于是
	$(A-\lambda_iI)^k\vect{c}_0=\vect{c}_k$，并且
	$(A-\lambda_iI)^{n_i}\vect{c}_0=\vect{0}$。反之，若这些递推关系成立，
	直接代回即可验证该函数是解。为得到非零解，必须有
	$\vect{c}_0\neq\vect{0}$。
\end{proof}

\begin{theorem}\label{thm:重特征值常系数方程组的基本解组}
	设 $A$ 有互不相同的特征值
	$\lambda_1,\lambda_2,\ldots,\lambda_s$，其重数分别为
	$n_1,n_2,\ldots,n_s$，且 $n_1+n_2+\cdots+n_s=n$。
	则齐次线性方程组有基本解组
	\[
		\begin{gathered}
		\mathrm{e}^{\lambda_1t}\vect{P}^{(1)}_1(t),
		\ldots,
		\mathrm{e}^{\lambda_1t}\vect{P}^{(1)}_{n_1}(t),\\
		\mathrm{e}^{\lambda_2t}\vect{P}^{(2)}_1(t),
		\ldots,
		\mathrm{e}^{\lambda_2t}\vect{P}^{(2)}_{n_2}(t),\\
		\hspace{3.5cm}\vdots\\
		\mathrm{e}^{\lambda_st}\vect{P}^{(s)}_1(t),
		\ldots,
		\mathrm{e}^{\lambda_st}\vect{P}^{(s)}_{n_s}(t).
		\end{gathered}
	\]
	这里，对每个 $i=1,2,\ldots,s$，
	\[
		\vect{P}^{(i)}_j(t)
		=\vect{c}^{(i)}_{j0}
		+\dfrac{t}{1!}\vect{c}^{(i)}_{j1}
		+\cdots
		+\dfrac{t^{n_i-1}}{(n_i-1)!}
		\vect{c}^{(i)}_{j(n_i-1)},
		\qquad j=1,2,\ldots,n_i,
	\]
	且零次项向量
	$\vect{c}^{(i)}_{10},\ldots,\vect{c}^{(i)}_{n_i0}$
	构成广义特征子空间 $V_i$ 的一组基，其余系数由递推关系
	\[
		\vect{c}^{(i)}_{jk}
		=(A-\lambda_iI)^k\vect{c}^{(i)}_{j0}
	\]
	确定。当 $A$ 为实矩阵时，可将共轭复解的实部与虚部分别取出，得到实基本解组。
\end{theorem}

\begin{proof}
	由引理~\ref{lem:向量多项式解的递推条件}，上述 $n$ 个函数都是方程组的解。
	令它们组成解矩阵 $X(t)$，则 $X(0)$ 的列向量恰为各子空间 $V_i$ 所选的基。
	由引理~\ref{lem:广义特征子空间分解}，这些列向量合在一起构成
	$\C^n$ 的一组基，故 $\det X(0)\neq0$。因此上述 $n$ 个解线性无关，
	构成方程组的基本解组。
\end{proof}
